\section{Threats to Validity}

\parabf{Internal.} 
One threat to validity comes from the data leakage of ground truth developer patches in \swebenchlite being part of the training data for \gptfouro. 
Since \gptfouro is a closed-source model, we do not have access to the training data.
Meanwhile, we note here that prior work almost exclusively used similar closed-source \llm{s} (e.g., \gptfouro, \gptfour, \claude-3.5, etc), and our approach can outperform all existing open-source solutions with same models.
Furthermore, the authors of \swebench~\cite{swebench} compared the resolve rate of issues collected before and after the knowledge cutoff date of \gptfour, and did not find any significant difference.  
To completely address this threat, we would need to retrain \gptfouro from scratch which would be infeasible for an academic project.

\parabf{External.}
One main external threat comes from our evaluation dataset of \swebenchlite.
While the performance of \tech might not generalize to other datasets, \swebenchlite is by far the most popular evaluation dataset which contains a diverse range of problems.
In addition, \emph{OpenAI has independently performed an extensive evaluation of \tech} and other open-source solutions on \swebenchlite, \swebench, and their new \swebenchverified benchmark, further confirming that \tech outperforms all other open-source agents~\cite{openaiverified}. Moreover, on Sept. 12th 2024, OpenAI just released the new o1 model family and also adopted Agentless as the top approach to showcase their performance on \swebench~\cite{openaioonesystemcard}.
In the future, we plan to further address this threat by evaluating \tech on other benchmarks.

\section{Conclusion}

We propose \tech -- an \emph{agentless} approach to automatically tackle software development problems. 
\tech uses a simple three phase approach of localization, repair, and patch validation. 
Compared to prior agent-based approaches, \tech deliberately disallows the \llm{ for autonomous tool usage or planning}. 
Our evaluation on the popular \swebenchlite benchmark demonstrates that \tech can achieve the highest performance compared with other open-source techniques while at the same time minimizing the cost. 
Furthermore, we perform a detailed classification of problems in \swebenchlite to not only offer new insights but to construct a more rigorous benchmark of \swebenchlitefiltered{} after removing problematic problems.

\section*{Acknowledgments}

We thank Jiawei Liu and Yuxiang Wei for providing some of the{ resources} used to run the experiments. One of the authors would like to thank Jun Yang for generously gifting his old bike\footnote{Sadly the bike is currently broken.} which allowed the author to travel faster and thus increasing research speed. 


